\documentclass{article}
\usepackage{booktabs, graphicx, hyperref, fontspec}
\usepackage{sectsty}
\allsectionsfont{\sffamily}
\usepackage[margin=1in]{geometry}
\hypersetup{
  colorlinks = true,
  urlcolor = cyan,
 }
 \providecommand{\tightlist}{%
  \setlength{\itemsep}{0pt}\setlength{\parskip}{0pt}}
\newcommand*{\authorfont}{\fontfamily{phv}\selectfont}
\usepackage[]{Fira Sans}

\begin{document}

\sffamily

\centerline{\Huge Microeconomic Analysis}

\vspace{3 mm}

\centerline{\large Dr.~Ryan Safner}
\vspace{2 mm}
\centerline{\large \href{http://microS22.classes.ryansafner.com}{microS22.classes.ryansafner.com}}

\vspace{5 mm}

\begin{tabular}{@{}p{3.5in}p{3.5in}}           
\textbf{Course}: ECON 306 Spring
2022  & \textbf{Email:}  \href{mailto:safner@hood.edu}{\nolinkurl{safner@hood.edu}}\\
\textbf{Room}: Rosenstock 212 & \textbf{Office:}  Rosenstock 110\\
\textbf{Meets}: MW 11:30 AM---12:55 PM & \textbf{Hours:}  MW 1:30
PM---2:30 PM \& by appt\\ 
\end{tabular}

\vspace{5 mm}

\hrule


\begin{quote}
``The curious task of economics is to demonstrate to men how little they
really know about what they imagine they can design.'' {--- F.A. Hayek
(1974 Economics Nobel Laureate)}
\end{quote}

\textbf{Microeconomic Analysis} is a collection of analytical tools for
modeling and understanding how decisions made by agents (individuals and
firms) in a context of market institutions lead to social coordination
and the creation of value. This toolkit is the analytical foundation for
exploring \emph{all} fields and applications of economics as a social
science, has many practical applications in business, finance, politics,
and insights into any realm of human endeavors. Economics, after all, is
the study of human action. At the intermediate level, these concepts are
made more rigorous through the use of mathematical models, but we will
first and foremost lay the theoretical foundations with economic
intuition. As such, this class assumes you have met the \textbf{required
prerequisite courses} by taking principles of microeconomics
(\textbf{ECON 206}).

I reserve the right to change any part of this syllabus and course, at
my discretion, with proper advance warning.

\hypertarget{course-format-and-covid}{%
\section{Course Format (and Covid)}\label{course-format-and-covid}}

As of Fall 2021, all students are expected to be on campus except those
with special approved exemptions. As such, this course will be taught
\textbf{in-person} and \textbf{synchronously} until or unless otherwise
announced.

You are expected to come to class except due to medical reasons or other
legitimate conflicts. Watching videos are not a substitute for attending
class.\footnote{On average, even for students who complete all
  assignments, those that do not regularly attend class suffer by a full
  letter grade,
  (\href{https://www.aeaweb.org/articles?id=10.1257/jep.7.3.167}{Levitt
  1993}).} Please see the \protect\hyperlink{attendance}{attendance
policy} for more.

In any event that we are unable to meet in person, I will hold class
meetings at the same day/time live on Zoom, and post all recorded
lectures via Panopto on Blackboard, and all assignments will be
submitted online (often via Blackboard).

\hypertarget{learning-during-a-global-pandemic}{%
\subsection{Learning During a Global
Pandemic}\label{learning-during-a-global-pandemic}}

While we have made some progress in returning to normal, this remains a
unique semester and a lot of things are still awful right now. None of
us signed up for this. None of us are really okay, we're all just
pretending for everyone else.

Many of you may be dealing with hardships at home and at work, and are
generally juggling many more problems than usual. Everyone's future
plans have been completely put on hold or cancelled to a large degree.

I am prioritizing us supporting each other as human beings during this
crazy era, and will try to use simple, accessible solutions that make
sense for the most people, and above all, to be flexible.

If you tell me you're having trouble, I will do whatever I can to help,
and not judge you or think less of you. I hope you will extend me the
same courtesy.

You never \emph{owe} me personal information about your health (mental
or physical). You are however always welcome to talk to me about things
that you're going through. If I can't help you, I usually know somebody
who can.

I want you to learn a lot from this course, but it is more important for
you to remain healthy, balanced, and grounded during this crisis.

\textbf{I reserve the right to change any part of this syllabus and
course, at my discretion, with proper advance warning.}

\hypertarget{course-objectives}{%
\section{Course objectives}\label{course-objectives}}

{By the end of this course,} you will:

\begin{enumerate}
\def\labelenumi{\arabic{enumi}.}
\tightlist
\item
  apply the models of microeconomics (constrained optimization and
  equilibrium) towards explaining real world behavior of individuals,
  firms, and governments
\item
  explore the effects of economic and political processes on market
  performance (competition, market prices, profits and losses, property
  rights, entrepreneurship, market power, market failures, public
  policy, government failures)
\item
  apply the economic way of thinking to real world issues in writing
\end{enumerate}

Given these objectives, this course fulfills all three of the learning
outcomes for
\href{https://www.hood.edu/academics/departments/george-b-delaplaine-jr-school-business/student-learning-outcomes}{the
George B. Delaplaine, Jr.~School of Business} Economics B.A. program:

\begin{itemize}
\tightlist
\item
  Use quantitative tools and techniques in the preparation,
  interpretation, analysis and presentation of data and information for
  problem solving and decision making {[}\ldots{]}
\item
  Apply economic reasoning and models to understand and analyze problems
  of public policy {[}\ldots{]}
\item
  Demonstrate effective oral and written communications skills for
  personal and professional success{[}\ldots{]}
\end{itemize}

{Fair warning:} \textbf{Economics is hard.} \emph{This, in particular,
may be of the hardest courses that you will take, primarily due to the
mathematical content.} Using the economic way of thinking is a skill, it
is literally retraining your brain to interpret and analyze the world in
a novel way, and is not something that can be memorized. I will do my
best to make this class intuitive and helpful, if not interesting. If at
any point you find yourself struggling in this course for any reason,
please come see me. Do not suffer in silence! Coming to see me for help
does not diminish my view of you, in fact I will hold you in
\emph{higher} regard for understanding your own needs and taking charge
of your own learning. There are also a some fantastic resources on
campus, such as the
\href{http://www.hood.edu/campus-services/academic-services/index.html}{Center
for Academic Achievement and Retention (CAAR)} and the
\href{http://www.hood.edu/library/}{Beneficial-Hodson Library}.

Free tutoring in many subjects is available for Hood students through
Thinking Storm professional tutoring and Hood peer tutoring. Hood peer
tutoring is located in the Tutoring Center adjacent to the Student
Success Center in the Beneficial Hodson Library and Learning Commons.
Appointments and online tutoring can be accessed through Blackboard and
located under tools/academic tutoring. Tutoring is available evenings
and weekends as well as during the day to help meet your needs. For more
information about tutoring and to see the specific subject tutoring
offered, please visit the student success website here or contact the
student success center at
\href{mailto:studentsuccess@hood.edu}{\nolinkurl{studentsuccess@hood.edu}}.

See my \href{/resources}{tips for success in this course}.

\hypertarget{required-course-materials}{%
\section{Required Course materials}\label{required-course-materials}}

You can find all course materials at my \textbf{dedicated website} for
this course:
\href{https://microS22.classes.ryansafner.com}{microS22.classes.ryansafner.com}.
Links to the website are posted on our Blackboard course page. Please
familiarize yourself with the website, see that it contains this
\href{syllabus/}{syllabus}, \href{resources/}{resources} to help you,
and our \href{schedule/}{schedule}. On the schedule page, you can find
each module with its own class page (\textbf{start there!}) along with
all related readings, lecture slides, practice problems, and
assignments.

My lecture slides will be shared with you, and serve as your primary
resource, but our main ``textbook'' below is \textbf{recommended} as the
next best resource and will be available from the campus bookstore. I
will discuss more about textbooks and materials in the first module.

\hypertarget{books}{%
\subsection{Books}\label{books}}

\begin{enumerate}
\def\labelenumi{\arabic{enumi}.}
\tightlist
\item
  Goolsbee, Austan, Steven Levitt, and Chad Syverson, 2012,
  \emph{Microeconomics} 2nd ed, USA: Worth Publishers, ISBN:
  978-1464187025
\end{enumerate}

You are welcome to purchase the book by other means (e.g.~Amazon,
half.com, etc). I have no financial stake in requiring you to purchase
this book. You are welcome to use previous version of the book, but
carefully verify the reading assignments, as the material may be
different across versions.

\hypertarget{assignments-and-grades}{%
\section{Assignments and Grades}\label{assignments-and-grades}}

Your final course grade is the weighted average of the following
assignments. You can find general descriptions for all the assignments
on the
\href{https://microS22.classes.ryansafner.com/assignments/}{assignments
page} and more specific information and examples on each assignment's
page on the
\href{https://microS22.classes.ryansafner.com/schedule/}{schedule page}.

\begin{tabular}{l|l|l}
\hline
Frequency & Assignment & Weight\\
\hline
1 & Opinion-Editorial & 20\%\\
\hline
5 & Homeworks & 20\% (using average HW grade)\\
\hline
3 & Exams & 20\% (each)\\
\hline
\end{tabular}

Each assignment is graded on a 100 point scale. Letter-grade equivalents
are based on the following scale:

\begin{table}
\centering
\begin{tabular}{l|c|l|c}
\hline
Grade & Range & Grade & Range\\
\hline
A & 93–100\% & C & 73–76\%\\
\hline
A− & 90–92\% & C− & 70–72\%\\
\hline
B+ & 87–89\% & D+ & 67–69\%\\
\hline
B & 83–86\% & D & 63–66\%\\
\hline
B− & 80–82\% & D− & 60–62\%\\
\hline
C+ & 77–79\% & F & < 60\%\\
\hline
\end{tabular}
\end{table}

See also my
\href{https://ryansafner.shinyapps.io/306_grade_calculator/}{
\texttt{Grade\ Calculator}} app where you can calculate your overall
grade using existing assignment grades and forecast ``what if''
scenarios.

These grades are firm cutoffs, but I do of course round upwards
\((\geq\) 0.5) for final grades. A necessary reminder, as an academic, I
am not in the business of \emph{giving} out grades, I merely report the
grade that you \emph{earn}. I will not alter your grade unless you
provide a reasonable argument that I am in error (which does happen from
time to time).

\textbf{No extra credit is available}

\hypertarget{policies-and-expectations}{%
\section{Policies and Expectations}\label{policies-and-expectations}}

This syllabus is a contract between you, the student, and me, your
instructor. It has been carefully and deliberately thought out. (A
syllabus can and will be used as a legal document for disputes tried at
a court of law. Ask me how I know.), and I will uphold my end of the
agreement and expect you to uphold yours.

In the language of game theory, this syllabus is my commitment device. I
am a very understanding person, and I know that exceptions to rules
often need to be made for students. However, to be \emph{fair} to
\emph{all} students the syllabus artificially constrains my ability to
make exceptions at a whim for anyone. This prevents clever students from
exploiting my congenial personality at everyone else's expense. Please
read and familiarize yourself with the course policies and expectations
of you. Chances are, if you have a question, it is answered herein.

\hypertarget{attendence}{%
\subsection{Attendence}\label{attendence}}

Your day-to-day classroom attendance is not graded. My philosophy is
that you are all adults and must take ownership of your own learning or
else you will not succeed. Some assignments may require in-class
participation for credit, and an (unexcused) absence may be detrimental
to your grade. Attending class is one of the strongest predictors of
success.

However, as required under Hood College's ``Promise of Fall Plan,'' (Ch.
2-C) \textbf{your classroom attendance will be recorded at every class
meeting.} This is primarily to facilitate contact tracing.

If you know you will be absent, you are not \emph{required} to let me
know, but it is polite to give notice (Note if I do not reply to an
email of yours letting me know, I am probably busy but will still see it
and appreciate your email). Your absence will be noted and recorded for
the purposes stated above. If, however, we have an assignment due in
class, you \emph{must} notify me ahead of time in order to make
alternate arrangements to still receive credit. Hasty ex-post attempts
to notify me will generate little sympathy.

\hypertarget{late-assignments}{%
\subsection{Late Assignments}\label{late-assignments}}

I will accept late assignments, but will subtract a specified amount of
points as a penalty. Even if it is the last week of the semester, I
encourage you to turn in late work: some points are better than no
points!

\textbf{Homeworks}: If you turn in a homework after it is due but before
it is graded or the answer key posted, I generally will not take off any
points. However, \textbf{if you turn in a homework \emph{after} the
answer key is posted, I will automatically deduct 20 points (so the
maximum grade you can earn on it is an 80).}

\textbf{Exams}: If you know that you will be unable to complete an
\emph{in-class exam} as scheduled for a legitimate reason, please notify
me at least \emph{one week} in advance, and we will schedule a make-up
exam date. Failure to do so, including desperate attempts to make
arrangements only \emph{after} the exam will result in a grade of 0 and
little sympathy.

\textbf{Op-eds}: Starting at the deadline, I will take off 1 point for
every hour that your Op-ed is late.

I reserve the right to re-weight assignments for students whom I believe
are legitimately unable to complete a particular assignment.

\hypertarget{grading}{%
\subsection{Grading}\label{grading}}

I will try my best to post grades on Blackboard's Grading Center and
return graded assignments to you within about one week of you turning
them in. There will be exceptions. Where applicable, I will post answer
keys once I know most homeworks are turned in (see Late Assignments
above for penalties). Blackboard's Grading Center is the place to look
for your most up-to-date grades. See also my
\href{https://ryansafner.shinyapps.io/306_grade_calculator/}{
\texttt{Grade\ Calculator}} app where you can calculate your overall
grade using existing assignment grades and forecast ``what if''
scenarios.

\hypertarget{communication-email-slack-and-virtual-office-hours}{%
\subsection{Communication: Email, Slack, and Virtual Office
Hours}\label{communication-email-slack-and-virtual-office-hours}}

Students must regularly monitor their \textbf{Hood email accounts} to
receive important college information, including messages related to
this class. Email through the Blackboard system is my main method of
communicating announcements and deadlines regarding your assignments.
\textbf{Please do not reply to any automated Blackboard emails - I may
not recieve it!}. My Hood email (\texttt{safner@hood.edu}) is the best
means of contacting me. I will do my best to respond within 24 hours. If
I do not reply within 48 hours, do not take it personally, and
\emph{feel free to send a follow up email} in the very likely event that
I genuinely did not see your original message.

Our \href{https://hoodcollegeeconomics.slack.com}{slack channel} is
available to all students and faculty in Economics and Business. I have
invited all of my classes and advisees. It will not be extended to
non-Business/Economics students or faculty. All users must use their
\textbf{hood emails} and \textbf{true first and last names}. Each course
has its own channel, exclusive for verified students in the course, and
myself, by my invite only. As a third party platform, you agree to its
Terms of Service. I have created this space as a way to stay connected,
to help one another, and to foster community. Behaviors such as posting
inappropriate content, harassing others, or engaging in academic
dishonesty, to be determined solely at my discretion, will result in one
warning, the content will be deleted, and subsequent behavior will
result in a ban.

In addition to in-person office hours, you can also make an appointment
for \textbf{``office hours''} on Zoom. You can join in with video,
audio, and/or chat, whichever you feel comfortable with. Of course, if
you are not available during those times, we can schedule our own time
if you prefer this method over email or Slack. If you want to go over
material from class, please have \emph{specific} questions you want help
with. I am not in the business of giving private lectures (particularly
if you missed class without a valid excuse).

Watch the excellent and accurate video explaining office hours
\href{https://gamef21.classes.ryansafner.com/syllabus/\#communication-email-slack-and-virtual-office-hours}{on
the website}.

\hypertarget{netiquette}{%
\subsection{Netiquette}\label{netiquette}}

When using Zoom and Slack, please follow appropriate internet etiquette
(``Netiquette''). Written communications, like blog posts or use of the
Zoom chat, lacks important nonverbal cues (such as body language, tone
of voice, sarcasm, etc).

Above all else, please respect one another and think/reread carefully
about how others may see your post before you submit a comment. You are
expected to disagree and have different opinions, this is inherently
valuable in a discussion. Please be civil and constructive in responding
to others' comments: writing \emph{``have you considered `X'?''} is a
lot more helpful to all involved than just writing \emph{``well you're
just wrong.''}

Posting content that is wilfully incindiary, illegal, or that
constitutes academic dishonesty (such as plagarism) will automatically
earn a grade of 0 and may be elevated to other authorities on campus.

When using the chat function on Zoom or public Slack channels, please
treat it as official course communications, even though I may not be
grading it. It may be a quick and informal tool - don't feel you need to
worry about spelling or perfect grammar - but please try to avoid
\emph{too} informal ``text-speak'' (i.e.~say ``That's good for you''
instead of ``thas good 4 u'').

\hypertarget{privacy}{%
\subsection{Privacy}\label{privacy}}

\href{https://www.execvision.io/blog/maryland-call-recording-laws/}{Maryland
law}
\href{https://law.justia.com/codes/maryland/2005/gcj/10-402.html}{requires}
all parties consent for a conversation or meeting to be recorded. If you
join in, and certainly if you participate, \textbf{you are consenting to
be recorded.} However, as described below, videos are \emph{not
accessible} beyond our class.

Live lectures are recorded on Zoom and posted to Blackboard via Panopto,
a secure course management system for video. Among other nice features
(such as multiple video screens, close captioning, and time-stamped
search functions!), Panopto is authenticated via your Blackboard
credentials, ensuring that \emph{our course videos are not accessible to
the open internet.}

For the privacy of your peers, and to foster an environment of trust and
academic freedom to explore ideas, \textbf{do not record our course
lectures or discussions.} You are already getting my official copies.

The
\href{https://www2.ed.gov/policy/gen/guid/fpco/ferpa/index.html}{Family
Educational Rights and Privacy Act} prevents me from disclosing or
discussing any student information, including grades and records about
student performance. If the student is at least 18 years of age,
\emph{parents (or spouses) do not have a right to obtain this
information}, except with consent by the student.

Many of you may be tuning in remotely, living with parents, and may have
occasional interruptions due to sharing a space. This is normal and
fine, but know that I will protect your privacy and not discuss your
performance when parents (or anyone other than you, for that matter) are
present, without your explicit consent.

\hypertarget{enrollment}{%
\subsection{Enrollment}\label{enrollment}}

Students are responsible for verifying their enrollment in this class.
The last day to add or drop this class with no penalty is
\textbf{Wednesday, September 1}. Be aware of
\href{https://www.hood.edu/offices-services/registrars-office/academic-calendar}{important
dates}.

\hypertarget{honor-code}{%
\subsection{Honor Code}\label{honor-code}}

Hood College has an Academic Honor Code which requires all members of
this community to maintain the highest standards of academic honesty and
integrity. Cheating, plagiarism, lying, and stealing are all prohibited.
All violations of the Honor Code are taken seriously, will be reported
to appropriate authority, and may result in severe penalties, including
expulsion from the college. See
\href{http://hood.smartcatalogiq.com/en/2016-2017/Catalog/The-Spirit-of-Hood/The-Academic-Honor-Code-and-Code-of-Conduct}{here}
for more detailed information.

\hypertarget{van-halen-and-mms}{%
\subsection{Van Halen and M\&Ms}\label{van-halen-and-mms}}

When you have completed reading the syllabus, email me a picture of the
band Van Halen and a picture of a bowl of M\&Ms.~If you do this
\emph{before} the date of the first exam, you will get bonus points on
the exam. If 75-100\% of the class does this, you each get 2 points. If
50-75\% of the class does this, you each get 4 points. If 25-50\% of the
class does this, you each get 6 points. If 0-25\% of the class does
this, you each get 8 points. Yes, you read this correctly.

\hypertarget{accessibility-equity-and-accommodations}{%
\subsection{Accessibility, Equity, and
Accommodations}\label{accessibility-equity-and-accommodations}}

College courses can, and should, be challenging and bring you out of
your comfort zone in a safe and equitable environment. If, however, you
feel at any point in the semester that certain assignments or aspects of
the course will be disproportionately uncomfortable or burdensome for
you due to any factor beyond your control, please come see me or email
me. I am a very understanding person and am happy to work out a solution
together. I reserve the right to modify and reweight assignments at my
sole discretion for students that I belive would legitimately be at a
disadvantage, through no fault of their own, to complete them as
described.

If you are unable to afford required textbooks or other resources for
any reason, come see me and we can find a solution that works for you.

This course is intended to be accessible for all students, including
those with mental, physical, or cognitive disabilities, illness,
injuries, impairments, or any other condition that tends to negatively
affect one's equal access to education. If at any point in the term, you
find yourself not able to fully access the space, content, and
experience of this course, you are welcome to contact me to discuss your
specific needs. I also encourage you to contact the
\href{https://www.hood.edu/academics/josephine-steiner-center-academic-achievement-retention/accessibility-services}{Office
of Accessibility Services} (301-696-3421). If you have a diagnosis or
history of accommodations in high school or previous postsecondary
institutions, Accessibility Services can help you document your needs
and create an accommodation plan. By making a plan through Accessibility
Services, you can ensure appropriate accommodations without disclosing
your condition or diagnosis to course instructors.

\hypertarget{acbsp-information}{%
\subsection{ACBSP Information}\label{acbsp-information}}

\begin{figure}[h!]
\centering
\includegraphics[]{../images/acbsp.png}
\end{figure}

Hood College is an accredited member of the Accreditation Council for
Business Schools and Programs (ACBSP), an organization devoted to
enhancing business education. In receiving and maintaining this
accreditation, the faculty has made a commitment to the continuous
improvement, innovation, and scholarship of the Department of Economics
and Management. For you, this means that your educational experience
undergoes ongoing validation to ensure it meets the most rigorous
international standards of business education. Only a select group of
institutions have received this status and it is an attribute of Hood
for which you should take great pride.

Pragmatically, our accreditation means that we will engage in the
ongoing use of measures, both quantitative and qualitative in nature, to
assess the performance of our students and the program. We ask that you
take very seriously the surveys and other measurement devices we will
use -- your best work and honest response will help us best assess and
improve our program.

\hypertarget{tentative-schedule}{%
\section{Tentative Schedule}\label{tentative-schedule}}

\textbf{You can find a full schedule} with much more details, including
the readings, appendices, and other further resources for each class
meeting on the
\href{https://microS22.classes.ryansafner.com/schedule/}{schedule page}.

\end{document}